\documentclass[10pt]{amsart}
\usepackage[a4paper,margin=23mm]{geometry}
\usepackage{amsmath,amssymb,amsthm,mathtools,hyperref,graphicx,xspace,letltxmacro,xpatch}
\usepackage[english]{babel}
\usepackage{xfrac,xcolor}

\usepackage[scr=rsfs,cal=euler]{mathalfa}
\makeatletter
\makeatother

\multlinegap0pt
\newcommand{\Psfrac}[2]{(\sfrac{#1}{#2})}
\newcommand{\psfrac}[2]{\sfrac{(#1)}{#2}}
\newcommand{\spfrac}[2]{\sfrac{#1}{(#2)}}
\newcommand{\pspfrac}[2]{\sfrac{(#1)}{(#2)}}
\newcommand{\sbullet}{{\scriptscriptstyle\bullet}}
\newcommand{\cbullet}{{\raisebox{1pt}{$\sbullet$}}}

\let\epsilon\varepsilon
\let\<\langle
\let\>\rangle

\DeclareFontFamily{U}{mathx}{\hyphenchar\font45}
\DeclareFontShape{U}{mathx}{m}{n}{
<5> <6> <7> <8> <9> <10>
<10.95> <12> <14.4> <17.28> <20.74> <24.88>
mathx10
}{}
\DeclareSymbolFont{mathx}{U}{mathx}{m}{n}
\DeclareMathAccent{\widecheck} {0}{mathx}{"71}

\usepackage{mathabx}
\usepackage{esint}
\usepackage{bm}

\usepackage{tikz}
\usetikzlibrary{intersections,decorations.markings}
\usepackage{pgfplots}
\usetikzlibrary{pgfplots.fillbetween}

\newcommand{\R}{\mathbb{R}}
\newcommand{\Z}{\mathbb{Z}}
\newcommand{\CC}{\mathbb{C}}\let\C\CC

\newcommand{\bv}{\bm{v}}
\let\eps\varepsilon
\newcommand{\cS}{\mathcal{S}}
\newcommand{\dd}{\,d}
\newcommand{\dmin}{d_{\min}}
\newcommand{\loc}{\mathrm{loc}}
\newcommand{\rY}{\mathrm{Y}}

\DeclareMathOperator{\Bog}{Bog}
\DeclareMathOperator{\diam}{diam}
\DeclareMathOperator{\divv}{div}\let\div\divv
\DeclareMathOperator{\Div}{Div}
\DeclareMathOperator{\dist}{dist}
\DeclareMathOperator{\Id}{Id}
\DeclareMathOperator{\supp}{supp}
\DeclareMathOperator{\Y}{Y}

\newtheorem{thm}{Theorem}[section]
\newtheorem{cor}[thm]{Corollary}
\newtheorem{Lemma}[thm]{Lemma}
\newtheorem{Proposition}[thm]{Proposition}
\theoremstyle{definition}
\newtheorem{Definition}[thm]{Definition}
\theoremstyle{remark}
\newtheorem{remark}[thm]{Remark}
\newtheorem{op}[thm]{Open Problem}

\let\dpl\displaystyle
\providecommand{\skpt}{\par\smallskip\noindent}
\providecommand{\pointir}{.\ }
\providecommand{\eg}{e.g.}\providecommand{\ie}{i.e.}\providecommand{\etc}{etc.}\providecommand{\resp}{resp.}\providecommand{\cf}{cf.}
\providecommand{\mainmatter}{}\providecommand{\backmatter}{}
\let\Subsection\subsection\let\Subsubsection\subsubsection
\AtBeginDocument{\gdef\shorttitle{Motion of several slender rigid filaments}}

\let\oldbigcup\bigcup

\pgfdeclarelayer{pre main}
\pgfsetlayers{pre main,main}
